\documentclass[11pt]{article}
\usepackage[margin=1in]{geometry}
\usepackage[T1]{fontenc}
\usepackage{lmodern}
\usepackage{microtype}
\usepackage{xcolor}
\usepackage{hyperref}
\usepackage{booktabs}
\usepackage{enumitem}
\usepackage{listings}
\usepackage{tcolorbox}
\usepackage{parskip}

\hypersetup{
  colorlinks=true,
  linkcolor=blue!60!black,
  urlcolor=blue!60!black
}

\definecolor{codebg}{RGB}{245,245,245}
\definecolor{codecomment}{RGB}{90,100,90}
\definecolor{codekeyword}{RGB}{0,70,140}
\definecolor{codestring}{RGB}{150,60,20}

\lstdefinestyle{pythonstyle}{
  language=Python,
  backgroundcolor=\color{codebg},
  basicstyle=\ttfamily\small,
  keywordstyle=\color{codekeyword}\bfseries,
  commentstyle=\color{codecomment}\itshape,
  stringstyle=\color{codestring},
  showstringspaces=false,
  breaklines=true,
  frame=single,
  rulecolor=\color{black!20},
  columns=fullflexible,
  keepspaces=true
}
\lstdefinestyle{shellstyle}{
  language=bash,
  backgroundcolor=\color{codebg},
  basicstyle=\ttfamily\small,
  showstringspaces=false,
  breaklines=true,
  frame=single,
  rulecolor=\color{black!20},
  columns=fullflexible,
  keepspaces=true
}
\lstset{style=pythonstyle}

\newtcolorbox{keyidea}{
  colback=blue!5, colframe=blue!50!black,
  title=Key Idea, fonttitle=\bfseries
}
\newtcolorbox{tipbox}{
  colback=green!5, colframe=green!45!black,
  title=Tip, fonttitle=\bfseries
}
\newtcolorbox{warningbox}{
  colback=orange!7, colframe=orange!60!black,
  title=Common Pitfall, fonttitle=\bfseries
}

\title{
  \textbf{Codebase Change Summary --- Week 3}\\[4pt]
  \large Equivalence Partitioning, Boundary-Value Analysis, and Negative Testing\\[2pt]
  \normalsize SE 413/513 --- Software Engineering
}
\author{Demo Testing Suite}
\date{}

\begin{document}
\maketitle

\begin{center}
\textit{What changed in \texttt{demo-testing-suite/} since Week 2, and why}
\end{center}

\vspace{0.3cm}

% ============================================================
\section{Purpose of This Document}
% ============================================================
The \texttt{demo-testing-suite/} codebase is a single, shared project that
grows one lecture at a time rather than being rebuilt from scratch each
week. This document summarizes exactly what was added or changed for
Week 3 and, more importantly, explains \emph{why} each change exists ---
what testing concept from the Week 3 reading it is meant to demonstrate.
Read this alongside the reading and \texttt{demo-testing-suite/README.md}.

\begin{keyidea}
Nothing from Weeks 1--2 was removed or rewritten. Everything below is
\emph{additive} --- the running examples you already saw
(\texttt{calculate\_late\_fee}, the shopping cart, discounts) are all
still there, unchanged, and still behave the same way.
\end{keyidea}

% ============================================================
\section{What Changed At a Glance}
% ============================================================
\begin{center}
\begin{tabular}{@{}p{5.4cm}p{1.6cm}p{7.4cm}@{}}
\toprule
\textbf{File} & \textbf{Status} & \textbf{Purpose} \\
\midrule
\texttt{app/pricing/shipping.py} & New & The Week 3 worked example: a
weight-based shipping-cost function with equivalence classes, boundaries,
and a deliberate defect for you to find. \\
\texttt{tests/pricing/test\_shipping.py} & New & A parametrized
\texttt{pytest} suite encoding equivalence partitioning, boundary-value
analysis, and negative testing against that function. \\
\texttt{docs/DEMO\_GUIDE\_week03\_ep\_bva\_negative\_testing.md} & New &
A section-by-section lecture companion (primarily for in-class use). \\
\texttt{README.md} & Modified & Updated project layout description and
the expected \texttt{pytest -v} output, now that Week 3's files exist. \\
\bottomrule
\end{tabular}
\end{center}

% ============================================================
\section{New App Code: \texttt{determine\_shipping\_cost()}}
% ============================================================
\textbf{Location:} \texttt{app/pricing/shipping.py}

This is the centerpiece example for this lecture. Its specification (in
the file's own docstring) is:

\begin{quote}
Given a package weight in kilograms, return the shipping cost: \$5.00 for
$0 < \text{weight} \leq 1$\,kg, \$10.00 for $1 < \text{weight} \leq
5$\,kg, and \$20.00 for $5 < \text{weight} \leq 20$\,kg. Weights of
0\,kg or less, or above 20\,kg, are invalid and should raise a
\texttt{ValueError}.
\end{quote}

This specification is deliberately identical in shape to the
\texttt{determine\_shipping\_cost()} worked example walked through
step by step in the Week 3 reading --- so you can use the reading's
equivalence-class table and boundary table directly against this real
file instead of against a hypothetical one.

\begin{warningbox}
\textbf{This implementation is not guaranteed to be correct.} It
contains at least one defect placed there on purpose, at one of the
specification's boundaries. Finding it is meant to be an exercise in
applying boundary-value analysis yourself --- work out the boundary
values from the specification first, predict the expected cost at each
one, \emph{then} check what the function actually returns. This document
intentionally does not name the exact defect; treat that as part of the
exercise.
\end{warningbox}

There is also a second, separate issue in how the function responds to
input of the wrong \emph{type} entirely (not just the wrong numeric
range) --- worth investigating once you've finished with the numeric
boundaries.

% ============================================================
\section{New Tests: \texttt{test\_shipping.py}}
% ============================================================
\textbf{Location:} \texttt{tests/pricing/test\_shipping.py}

This file encodes the three techniques from the reading as one
parametrized \texttt{pytest} suite, in the same order the reading
builds them:

\begin{itemize}
\item \textbf{Equivalence partitioning:} one representative value from
each valid class (below the first boundary, between each pair of
boundaries, above the last valid boundary).
\item \textbf{Boundary-value analysis:} the exact edge of every class in
the specification, tested together with the equivalence-partitioning
cases in one parametrized test function.
\item \textbf{Negative testing:} two separate test functions --- one
supplying out-of-range numeric weights and asserting a
\texttt{ValueError}, and one supplying values of the wrong type
entirely (\texttt{None}, a string, a list, a dictionary) and asserting
some kind of clear error.
\end{itemize}

\begin{tipbox}
Run the suite yourself before reading further:
\begin{lstlisting}[style=shellstyle]
cd demo-testing-suite
pytest tests/pricing/test_shipping.py -v
\end{lstlisting}
One case in the output is marked \texttt{XFAIL} (``expected failure'')
rather than \texttt{PASS} or \texttt{FAIL}. That is not a broken test ---
it is how this codebase documents a \emph{known}, not-yet-fixed defect,
exactly as \texttt{calculate\_late\_fee}'s negative-days bug was
documented back in Week 1. Use the \texttt{XFAIL} line as a hint about
\emph{where} the boundary defect lives, without reading the reason
string just yet if you want to work it out yourself first.
\end{tipbox}

% ============================================================
\section{New Documentation}
% ============================================================
\textbf{Location:} \texttt{docs/DEMO\_GUIDE\_week03\_ep\_bva\_negative\_testing.md}

This is a lecture-companion document, primarily written to guide what
gets shown and run during class, section by section. It is organized to
follow the Week 3 reading's own structure. You are welcome to read it,
but it is not required reading --- the reading itself and the exercise
above are the primary material.

% ============================================================
\section{Updated: \texttt{README.md}}
% ============================================================
The top-level \texttt{demo-testing-suite/README.md} was updated to
reflect the two new files above and to describe the suite's current
expected output. If you run the full suite from the project root:

\begin{lstlisting}[style=shellstyle]
cd demo-testing-suite
pip install -r requirements.txt
pytest -v
\end{lstlisting}

you should see a mix of passing tests, exactly two \texttt{XFAIL}
results, and exactly one \texttt{FAIL} (all three explained in the
README and the demo guide). If your output differs from that --- more
failures, fewer passes --- something is broken and worth investigating
or asking about; if it matches, everything is behaving as designed.

% ============================================================
\section{How This Connects to the Reading}
% ============================================================
\begin{keyidea}
The reading argues that equivalence partitioning gives you a small,
representative set of inputs; boundary-value analysis tells you which
values \emph{within} that set are most worth testing; and negative
testing checks that the software fails gracefully on input the
specification never anticipated. \texttt{shipping.py} and
\texttt{test\_shipping.py} exist so you can apply all three techniques
against a real function and a real, runnable test suite --- not just a
diagram on a slide.
\end{keyidea}

% ============================================================
\section{Suggested Next Steps}
% ============================================================
\begin{enumerate}
\item Read \texttt{app/pricing/shipping.py}'s specification (the
docstring), and independently derive its equivalence classes and
boundary values, the same way the reading does for its own example.
\item Predict the expected output at every boundary value \emph{before}
running anything.
\item Run \texttt{pytest tests/pricing/test\_shipping.py -v} and compare
the actual results to your predictions. Investigate any mismatch.
\item Try at least one additional negative test of your own --- an input
category not already covered in the file --- and reason about what the
function \emph{should} do with it, per the reading's guidance on
graceful failure.
\end{enumerate}

\end{document}
